\documentclass[aspectratio=169,11pt,french]{beamer}

\usepackage[T1]{fontenc}
\usepackage[utf8]{inputenc}
\usepackage[french]{babel}
\usepackage{lmodern}
\usepackage{graphicx}
\usepackage{tikz}
\usetikzlibrary{arrows.meta,positioning}

\definecolor{StagePurple}{HTML}{3434A8}
\definecolor{StageDark}{HTML}{25283A}
\definecolor{StagePale}{HTML}{F3F3FA}
\definecolor{StageGrey}{HTML}{626779}
\definecolor{StageTeal}{HTML}{178F8B}

\usetheme{default}
\setbeamercolor{normal text}{fg=StageDark,bg=white}
\setbeamercolor{structure}{fg=StagePurple}
\setbeamercolor{frametitle}{fg=white,bg=StagePurple}
\setbeamercolor{block title}{fg=white,bg=StagePurple}
\setbeamercolor{block body}{fg=StageDark,bg=StagePale}
\setbeamertemplate{navigation symbols}{}
\setbeamertemplate{itemize item}{\color{StagePurple}\small$\bullet$}
\setbeamersize{text margin left=0.8cm,text margin right=0.8cm}
\setlength{\parskip}{0.25em}

\setbeamerfont{frametitle}{series=\bfseries,size=\large}
\setbeamertemplate{frametitle}{%
  \nointerlineskip
  \begin{beamercolorbox}[wd=\paperwidth,ht=0.8cm,dp=0.3cm,leftskip=0.8cm]{frametitle}
    \usebeamerfont{frametitle}\insertframetitle
  \end{beamercolorbox}
}
\setbeamertemplate{footline}{%
  \leavevmode
  \hbox{%
    \begin{beamercolorbox}[wd=.77\paperwidth,ht=2.7ex,dp=1.1ex,leftskip=.8cm]{author in head/foot}
      \color{StageGrey}\scriptsize Mouhamed Cisse \quad | \quad Prodigious France
    \end{beamercolorbox}%
    \begin{beamercolorbox}[wd=.23\paperwidth,ht=2.7ex,dp=1.1ex,rightskip=.8cm plus1fil]{date in head/foot}
      \color{StageGrey}\scriptsize Spielberg \hfill \insertframenumber/\inserttotalframenumber
    \end{beamercolorbox}%
  }%
}

\tikzset{
  flow/.style={draw=StagePurple!55,fill=StagePale,rounded corners=2pt,
    align=center,text=StageDark,minimum height=0.9cm,inner sep=5pt},
  emphasis/.style={draw=StageTeal,fill=StagePale,rounded corners=2pt,
    align=center,text=StageDark,minimum height=0.9cm,inner sep=6pt},
  arrow/.style={-{Stealth[length=2mm]},draw=StageTeal,line width=0.9pt}
}

% Use the original report figure when available; otherwise draw its reported
% architecture directly so the slide remains complete without figure assets.
\newcommand{\reportfigure}[4]{%
  \IfFileExists{#1}{%
    \includegraphics[width=.96\linewidth,height=.76\textheight,keepaspectratio]{#1}%
  }{%
    #4
  }%
  \par\vspace{0.08cm}
  {\scriptsize\color{StageGrey}#3}
}

\newcommand{\kpi}[2]{%
  \begin{beamercolorbox}[wd=\linewidth,sep=.35cm,center,rounded=true]{block body}
    {\color{StagePurple}\fontsize{23}{26}\selectfont\bfseries #1}\\[-.1em]
    {\small\color{StageGrey}#2}
  \end{beamercolorbox}
}

\newcommand{\objectivecard}[3]{%
  \begin{beamercolorbox}[wd=\linewidth,ht=3.0cm,dp=.25cm,center,rounded=true]{block body}
    {\color{StagePurple}\fontsize{25}{27}\selectfont\bfseries #1}\\[.15cm]
    {\large\bfseries #2}\\[.25cm]
    {\small\color{StageGrey}#3}
  \end{beamercolorbox}
}

\newcommand{\architecturefallback}[1]{%
  \begin{tikzpicture}[node distance=.45cm,scale=.88,transform shape]
    #1
  \end{tikzpicture}%
}

\title{Spielberg}
\subtitle{Centraliser les données et orchestrer les usages IA}
\author{Mouhamed Cisse}
\institute{Prodigious France -- Publicis Groupe}
\date{11 mai -- 17 juillet 2026}

\begin{document}

% 0:20
\begin{frame}[plain]
  \vspace*{1.05cm}
  \begin{beamercolorbox}[wd=\paperwidth,sep=.55cm,leftskip=.8cm]{frametitle}
    {\color{white}\fontsize{30}{34}\selectfont\bfseries Spielberg}\\[.18cm]
    {\color{white}\Large Centraliser les données, orchestrer les usages IA}
  \end{beamercolorbox}
  \vspace{.8cm}
  {\LARGE\bfseries Mouhamed Cisse}\\[.35cm]
  {\large Soutenance de stage -- IA et Data Engineering}\\[.5cm]
  {\normalsize Prodigious France \quad | \quad Publicis Groupe \quad | \quad
   11 mai -- 17 juillet 2026}
  \vfill
  {\small\color{StageGrey}Tuteur entreprise : Rony Leborgne
    \hfill Encadrante académique : Juhyun Park}
\end{frame}

% 0:30 -- Conservé dans l'esprit de la première version appréciée.
\begin{frame}{Introduction et environnement de stage}
  \vspace{.35cm}
  \begin{columns}[T,onlytextwidth]
    \column{.48\textwidth}
      {\Large\bfseries Prodigious France}\\[.25cm]
      \begin{itemize}
        \item Entité de production et d'exécution créative de Publicis Groupe
        \item Création, production, post-production, digital et opérations
      \end{itemize}
    \column{.48\textwidth}
      {\Large\bfseries Mon stage}\\[.25cm]
      \begin{itemize}
        \item Département Développement, pôle Digital
        \item Poste : \textit{Assistant Lead Automation}
        \item Travail transversal avec les équipes staffing, finance, RH, IT et sécurité
      \end{itemize}
  \end{columns}
  \vfill
  \begin{beamercolorbox}[wd=\linewidth,sep=.4cm,rounded=true]{block body}
    \centering\large
    Un projet au croisement des besoins métier, de la donnée et du développement logiciel.
  \end{beamercolorbox}
\end{frame}

% 0:30
\begin{frame}{Contexte et problématique}
  \begin{center}
    {\large\bfseries Les données opérationnelles sont réparties entre plusieurs outils}
  \end{center}
  \vspace{.1cm}
  \begin{center}
  \begin{tikzpicture}[node distance=.8cm]
    \node[flow,text width=4.1cm,minimum height=2.0cm] (sources)
      {\textbf{Sources métier}\\[.1cm]\scriptsize
       Teamium : staffing\\
       Altair : projets, budgets, BF\\
       OneClick / rOar : RH, timesheets};
    \node[emphasis,text width=3.3cm,minimum height=2.0cm,right=.7cm of sources]
      (spielberg) {\Large\bfseries Spielberg\\[.15cm]\normalsize point d'accès\\unifié visé};
    \node[flow,text width=2.0cm,minimum height=1.4cm,right=.7cm of spielberg]
      (users) {Équipes métier};
    \draw[arrow] (sources) -- (spielberg);
    \draw[arrow] (spielberg.east) -- (users.west);
  \end{tikzpicture}
  \end{center}
  \vspace{.1cm}
  \begin{center}
    \small Recherche manuelle \quad | \quad Ressaisies \quad | \quad
    Visibilité fragmentée \quad | \quad Réconciliations difficiles
  \end{center}
\end{frame}

% 0:25
\begin{frame}{Objectif du stage}
  \vspace{.35cm}
  \begin{center}
    {\Large\bfseries Faire évoluer le POC Spielberg vers une plateforme plus exploitable}
  \end{center}
  \vspace{.75cm}
  \begin{columns}[T,onlytextwidth]
    \column{.31\textwidth}
      \objectivecard{01}{Structurer la donnée}{JSON vers SQLite et API métier}
    \column{.31\textwidth}
      \objectivecard{02}{Outiller les parcours}{Création projet, BDC et Kissflow}
    \column{.31\textwidth}
      \objectivecard{03}{Encadrer l'assistant}{Workers, garde-fous et sorties structurées}
  \end{columns}
\end{frame}

% 0:25
\begin{frame}{Méthodologie}
  \vspace{.55cm}
  \begin{center}
  \begin{tikzpicture}[node distance=.55cm]
    \node[flow,text width=2.6cm,minimum height=1.5cm] (a)
      {\large\bfseries 1. Cadrer\\[.15cm]\normalsize CDC et ateliers métier};
    \node[flow,text width=2.6cm,minimum height=1.5cm,right=of a] (b)
      {\large\bfseries 2. Structurer\\[.15cm]\normalsize Données et API};
    \node[flow,text width=2.6cm,minimum height=1.5cm,right=of b] (c)
      {\large\bfseries 3. Intégrer\\[.15cm]\normalsize Interfaces et agents};
    \node[flow,text width=2.6cm,minimum height=1.5cm,right=of c] (d)
      {\large\bfseries 4. Évaluer\\[.15cm]\normalsize Tests et LangSmith};
    \draw[arrow] (a) -- (b);
    \draw[arrow] (b) -- (c);
    \draw[arrow] (c) -- (d);
  \end{tikzpicture}
  \end{center}
  \vspace{.65cm}
  \begin{center}
    \large Ateliers métier réguliers \quad$\rightarrow$\quad développement incrémental
    \quad$\rightarrow$\quad tests des parcours
  \end{center}
  \vspace{.25cm}
  \begin{center}
    Séparer les décisions du LLM des validations et exécutions déterministes.
  \end{center}
\end{frame}

% 0:30
\begin{frame}{Contributions principales}
  \vspace{.3cm}
  \begin{itemize}
    \item Cadrage métier, contrat d'ingestion et préparation de la couche LLM
    \vspace{.35cm}
    \item Migration des données de démonstration JSON vers SQLite
    \vspace{.35cm}
    \item Backend FastAPI : schémas Pydantic, services et repositories
    \vspace{.35cm}
    \item Interfaces de création projet et de demande BDC ; intégration Kissflow
    \vspace{.35cm}
    \item Évolution de l'assistant jusqu'au Supervisor LangGraph, ajout du worker RAG
  \end{itemize}
\end{frame}

% 0:25 -- Vue d'ensemble reprise du rapport.
\begin{frame}{Architecture globale -- périmètre Spielberg}
  \centering
  \reportfigure{Figures/scope_spielberg.png}
    {Cadrage, pipeline ETL local, backend FastAPI, interfaces métier, Kissflow et couche IA.}
    {Architecture réalisée pendant le stage -- schéma de synthèse du rapport.}
    {\begin{tikzpicture}[node distance=.55cm,scale=.9,transform shape]
      \node[flow,text width=2.3cm,minimum height=1.15cm] (sources)
        {Sources\\métier};
      \node[flow,text width=2.3cm,minimum height=1.15cm,right=of sources] (etl)
        {Pipeline ETL\\et SQLite};
      \node[flow,text width=2.3cm,minimum height=1.15cm,right=of etl] (api)
        {Backend\\FastAPI};
      \node[emphasis,text width=2.4cm,minimum height=1.15cm,right=of api] (front)
        {Interfaces\\métier};
      \node[flow,text width=2.4cm,minimum height=1.15cm,below=.8cm of api] (kiss)
        {Workflow BDC\\Kissflow};
      \node[emphasis,text width=2.4cm,minimum height=1.15cm,below=.8cm of front] (ai)
        {Supervisor\\et workers IA};
      \draw[arrow] (sources) -- (etl);
      \draw[arrow] (etl) -- (api);
      \draw[arrow] (api) -- (front);
      \draw[arrow] (api) -- (kiss);
      \draw[arrow] (api) -- (ai);
    \end{tikzpicture}}
\end{frame}

% 0:25 -- Image dédiée : affichée en grand si le fichier du rapport est disponible.
\begin{frame}{Pipeline ETL -- des mocks vers SQLite}
  \centering
  \reportfigure{Figures/ETL_Schema.png}
    {Pipeline ETL des données Spielberg vers SQLite.}
    {Pipeline ETL des données Spielberg vers SQLite -- schéma du rapport.}
    {\begin{tikzpicture}[node distance=.65cm,scale=.9,transform shape]
      \node[flow,text width=2.5cm,minimum height=1.25cm] (source)
        {Teamium\\et données mock};
      \node[flow,text width=2.5cm,minimum height=1.25cm,right=of source] (load)
        {Chargement\\et normalisation};
      \node[flow,text width=2.5cm,minimum height=1.25cm,right=of load] (model)
        {Tables\\SQLite};
      \node[emphasis,text width=2.5cm,minimum height=1.25cm,right=of model] (use)
        {API et worker\\\texttt{data}};
      \draw[arrow] (source) -- (load);
      \draw[arrow] (load) -- (model);
      \draw[arrow] (model) -- (use);
      \node[below=.55cm of load,text=StageGrey,align=center]
        {Données nettoyées et structurées};
    \end{tikzpicture}}
\end{frame}

% 0:25
\begin{frame}{Backend FastAPI -- responsabilités séparées}
  \centering
  \reportfigure{Figures/shema_backend.png}
    {Architecture backend : routes, schémas, services, repositories, SQLite et intégrations.}
    {Architecture du backend FastAPI -- schéma du rapport.}
    {\begin{tikzpicture}[node distance=.55cm,scale=.82,transform shape]
      \node[emphasis,text width=2.35cm,minimum height=1cm] (front)
        {Frontend\\interfaces};
      \node[flow,text width=2.35cm,minimum height=1cm,right=of front] (routes)
        {Routes\\FastAPI};
      \node[flow,text width=2.35cm,minimum height=1cm,right=of routes] (schemas)
        {Schémas\\Pydantic};
      \node[flow,text width=2.35cm,minimum height=1cm,right=of schemas] (services)
        {Services\\métier};
      \node[flow,text width=2.35cm,minimum height=1cm,right=of services] (repo)
        {Repositories};
      \node[emphasis,text width=2.2cm,minimum height=.9cm,below=.75cm of repo] (db)
        {SQLite};
      \node[flow,text width=2.2cm,minimum height=.9cm,below=.75cm of services] (external)
        {Kissflow / agents};
      \draw[arrow] (front) -- (routes);
      \draw[arrow] (routes) -- (schemas);
      \draw[arrow] (schemas) -- (services);
      \draw[arrow] (services) -- (repo);
      \draw[arrow] (repo) -- (db);
      \draw[arrow] (services) -- (external);
    \end{tikzpicture}}
\end{frame}

% 0:25
\begin{frame}{Workflow métier -- demande BDC}
  \centering
  \reportfigure{Figures/workflow_creation_bdc.png}
    {Workflow applicatif de création et de soumission d'une demande BDC vers Kissflow.}
    {Workflow BDC -- schéma du rapport.}
    {\begin{tikzpicture}[node distance=.5cm,scale=.88,transform shape]
      \node[flow,text width=2.25cm,minimum height=1.1cm] (form)
        {Formulaire\\frontend};
      \node[flow,text width=2.25cm,minimum height=1.1cm,right=of form] (validate)
        {Validation\\Pydantic};
      \node[flow,text width=2.25cm,minimum height=1.1cm,right=of validate] (service)
        {Service et\\mapping Kissflow};
      \node[emphasis,text width=2.25cm,minimum height=1.1cm,right=of service] (api)
        {Kissflow\\create puis submit};
      \node[flow,text width=2.25cm,minimum height=1.1cm,right=of api] (db)
        {Trace locale\\SQLite};
      \draw[arrow] (form) -- (validate);
      \draw[arrow] (validate) -- (service);
      \draw[arrow] (service) -- (api);
      \draw[arrow] (api) -- (db);
      \node[below=.55cm of service,text=StageGrey,align=center]
        {Identifiant Spielberg \quad|\quad Statut \quad|\quad Identifiants Kissflow};
    \end{tikzpicture}}
\end{frame}

% 0:25
\begin{frame}{Workflow métier -- création de projet}
  \centering
  \reportfigure{Figures/workflow_creation_projet.png}
    {Workflow applicatif de création d'un projet dans Spielberg.}
    {Création projet -- wizard, validation, service métier et persistance SQLite.}
    {\architecturefallback{
      \node[flow,text width=2.25cm,minimum height=1.15cm] (wizard)
        {Wizard frontend\\4 étapes};
      \node[flow,text width=2.25cm,minimum height=1.15cm,right=of wizard] (api)
        {POST endpoint\\validation Pydantic};
      \node[flow,text width=2.25cm,minimum height=1.15cm,right=of api] (service)
        {Service de\\création projet};
      \node[emphasis,text width=2.25cm,minimum height=1.15cm,right=of service] (db)
        {Repository\\SQLite};
      \node[flow,text width=2.25cm,minimum height=1.15cm,right=of db] (views)
        {Liste, détail\\et dashboard};
      \draw[arrow] (wizard) -- (api);
      \draw[arrow] (api) -- (service);
      \draw[arrow] (service) -- (db);
      \draw[arrow] (db) -- (views);
    }}
\end{frame}

% 0:25
\begin{frame}{Assistant -- évolution de l'architecture}
  \vspace{.4cm}
  \begin{center}
  \begin{tikzpicture}[node distance=.7cm]
    \node[flow,text width=2.45cm,minimum height=1.7cm] (a)
      {\normalsize\bfseries Chat déterministe\\[.12cm]\small Trop rigide};
    \node[flow,text width=2.45cm,minimum height=1.7cm,right=of a] (b)
      {\normalsize\bfseries Agent LangChain\\[.12cm]\small Trop centralisé};
    \node[flow,text width=2.45cm,minimum height=1.7cm,right=of b] (c)
      {\normalsize\bfseries Router\\[.12cm]\small Contexte fragile};
    \node[emphasis,text width=2.45cm,minimum height=1.7cm,right=of c] (d)
      {\normalsize\bfseries Supervisor\\LangGraph\\[.12cm]\small Un interlocuteur};
    \draw[arrow] (a) -- (b);
    \draw[arrow] (b) -- (c);
    \draw[arrow] (c) -- (d);
  \end{tikzpicture}
  \end{center}
  \vspace{.7cm}
  \begin{center}
    {\large Le Supervisor garde le fil de la conversation ;
    les workers se spécialisent dans l'exécution.}
  \end{center}
\end{frame}

% 0:25
\begin{frame}{Première étape agentique -- agent unique}
  \centering
  \reportfigure{Figures/single_agent.png}
    {Un agent LangChain centralise le raisonnement, le prompt, la mémoire et les outils.}
    {Architecture du premier agent unique Spielberg.}
    {\architecturefallback{
      \node[flow,text width=2.3cm,minimum height=1.0cm] (user)
        {Question\\utilisateur};
      \node[emphasis,text width=3.0cm,minimum height=1.3cm,right=of user] (agent)
        {Agent unique\\LangChain + LLM};
      \node[flow,text width=2.5cm,minimum height=1.0cm,above right=.55cm and .75cm of agent] (prompt)
        {Prompt global\\et mémoire};
      \node[flow,text width=2.5cm,minimum height=1.0cm,below right=.55cm and .75cm of agent] (tools)
        {Tool SQL\\et données};
      \node[flow,text width=2.4cm,minimum height=1.0cm,right=1.0cm of agent] (result)
        {Réponse\\structurée};
      \draw[arrow] (user) -- (agent);
      \draw[arrow] (prompt) -- (agent);
      \draw[arrow] (agent) -- (tools);
      \draw[arrow] (tools) -- (result);
      \draw[arrow] (agent) -- (result);
    }}
\end{frame}

% 0:25
\begin{frame}{Première architecture multi-agent -- Router}
  \centering
  \reportfigure{Figures/schema_router.png}
    {Le Router classe la demande puis la délègue à un agent spécialisé.}
    {Architecture Router Spielberg : classification puis délégation.}
    {\architecturefallback{
      \node[flow,text width=2.2cm,minimum height=1.0cm] (user)
        {Demande};
      \node[emphasis,text width=2.8cm,minimum height=1.15cm,right=of user] (router)
        {Router\\classification};
      \node[flow,text width=2.15cm,minimum height=.9cm,above right=.55cm and .7cm of router] (data)
        {\texttt{data\_agent}};
      \node[flow,text width=2.15cm,minimum height=.9cm,right=of router] (project)
        {Création projet};
      \node[flow,text width=2.15cm,minimum height=.9cm,below right=.55cm and .7cm of router] (bdc)
        {\texttt{create\_bdc}};
      \node[flow,text width=2.0cm,minimum height=1.0cm,right=1.0cm of project] (answer)
        {Agent spécialisé\\et réponse};
      \draw[arrow] (user) -- (router);
      \draw[arrow] (router) -- (data);
      \draw[arrow] (router) -- (project);
      \draw[arrow] (router) -- (bdc);
      \draw[arrow] (data) -- (answer);
      \draw[arrow] (project) -- (answer);
      \draw[arrow] (bdc) -- (answer);
    }}
\end{frame}

% 0:20
\begin{frame}{Benchmark -- coût en tokens selon l'architecture}
  \centering
  \reportfigure{Figures/token_bench.png}
    {Comparaison du coût en tokens entre agent unique, Supervisor et swarm selon les domaines distracteurs.}
    {Comparaison du benchmark cité dans le rapport ; schéma qualitatif si la figure source est absente.}
    {\begin{tikzpicture}[x=1.0cm,y=.72cm]
      \draw[-{Stealth[length=2mm]},StageDark,line width=.8pt] (0,0) -- (9.8,0);
      \draw[-{Stealth[length=2mm]},StageDark,line width=.8pt] (0,0) -- (0,4.9);
      \draw[StageGrey!18] (0,1) -- (9.25,1);
      \draw[StageGrey!18] (0,2) -- (9.25,2);
      \draw[StageGrey!18] (0,3) -- (9.25,3);
      \node[rotate=90] at (-.55,2.45) {\small Tendance qualitative des tokens};
      \node[below] at (.65,-.08) {\scriptsize peu};
      \node[below] at (8.55,-.08) {\scriptsize nombreux};
      \node[below] at (4.8,-.5) {\small Domaines distracteurs};
      \draw[StagePurple,line width=2.2pt,smooth]
        plot coordinates {(.55,.55) (2.2,.78) (4.3,1.35) (6.3,2.35) (8.65,3.8)};
      \draw[StageTeal,line width=2.2pt,smooth]
        plot coordinates {(.55,.72) (2.2,.94) (4.3,1.35) (6.3,1.8) (8.65,2.3)};
      \draw[StageGrey,line width=2.2pt,smooth]
        plot coordinates {(.55,.9) (2.2,1.1) (4.3,1.3) (6.3,1.5) (8.65,1.65)};
      \node[anchor=west,text=StagePurple] at (8.72,3.8) {\scriptsize Agent unique};
      \node[anchor=west,text=StageTeal] at (8.72,2.3) {\scriptsize Supervisor};
      \node[anchor=west,text=StageGrey] at (8.72,1.65) {\scriptsize Swarm};
    \end{tikzpicture}}
\end{frame}

% 0:30 -- Schéma original du rapport en pleine diapositive.
\begin{frame}{Architecture Supervisor -- graphe LangGraph}
  \centering
  \reportfigure{Figures/supervisor_graph.png}
    {Graphe LangGraph du Supervisor Spielberg.}
    {Graphe LangGraph du Supervisor Spielberg -- schéma du rapport.}
    {\begin{tikzpicture}[node distance=.5cm and .65cm,scale=.94,transform shape]
      \node[flow,text width=3cm] (guard) {\texttt{input\_guard}};
      \node[emphasis,text width=3.5cm,below=of guard] (decision)
        {\scriptsize\texttt{supervisor\_decision}};
      \node[flow,text width=2.15cm,below left=.7cm and 1.1cm of decision] (data)
        {\texttt{data}};
      \node[flow,text width=2.15cm,below=.7cm of decision] (project)
        {\texttt{project}};
      \node[flow,text width=2.15cm,below right=.7cm and 1.1cm of decision] (bdc)
        {\texttt{bdc}};
      \node[flow,text width=2.15cm,right=.5cm of bdc] (rag) {\texttt{rag}};
      \node[emphasis,text width=4.1cm,below=.65cm of project] (final)
        {\scriptsize\texttt{finalize\_from\_worker}};
      \draw[arrow] (guard) -- (decision);
      \draw[arrow] (decision) -- (data);
      \draw[arrow] (decision) -- (project);
      \draw[arrow] (decision) -- (bdc);
      \draw[arrow] (decision) -- (rag);
      \draw[arrow] (data) -- (final);
      \draw[arrow] (project) -- (final);
      \draw[arrow] (bdc) -- (final);
      \draw[arrow] (rag) -- (final);
    \end{tikzpicture}}
\end{frame}

% 0:25 -- Schéma original du rapport en pleine diapositive.
\begin{frame}{Worker RAG -- recherche documentaire sourcée}
  \centering
  \reportfigure{Figures/schema_rag_worker.png}
    {Détection de l'intention, recherche vectorielle, sélection des sources et réponse.}
    {Pipeline RAG intégré au Supervisor -- schéma du rapport.}
    {\begin{tikzpicture}[node distance=.65cm,scale=.88,transform shape]
      \node[flow,text width=2.25cm,minimum height=1.05cm] (docs)
        {Documents\\PDF / TXT / DOCX};
      \node[flow,text width=2.25cm,minimum height=1.05cm,right=of docs] (chunks)
        {Extraction\\et découpage};
      \node[flow,text width=2.25cm,minimum height=1.05cm,right=of chunks] (embed)
        {Embeddings\\et indexation};
      \node[emphasis,text width=2.25cm,minimum height=1.05cm,right=of embed] (chroma)
        {Chroma\\recherche};
      \node[flow,text width=2.25cm,minimum height=1.05cm,right=of chroma] (answer)
        {Worker RAG\\réponse sourcée};
      \draw[arrow] (docs) -- (chunks);
      \draw[arrow] (chunks) -- (embed);
      \draw[arrow] (embed) -- (chroma);
      \draw[arrow] (chroma) -- (answer);
      \node[below=.5cm of chroma,text=StageGrey,align=center]
        {Sources affichées dans le frontend};
    \end{tikzpicture}}
\end{frame}

% 0:40
\begin{frame}{Évaluation LangSmith -- résultats}
  \vspace{.35cm}
  \begin{center}
    {\Large\bfseries 27 scénarios répétés deux fois \quad | \quad 54 exécutions}
  \end{center}
  \vspace{.5cm}
  \begin{columns}[T,onlytextwidth]
    \column{.32\textwidth}
      \kpi{54/54}{succès techniques}
    \column{.32\textwidth}
      \kpi{48/54}{score structurel\\88,9\,\%}
    \column{.32\textwidth}
      \kpi{45/54}{score textuel strict\\83,3\,\%}
  \end{columns}
  \vspace{.55cm}
  \begin{columns}[T,onlytextwidth]
    \column{.48\textwidth}
      \begin{beamercolorbox}[wd=\linewidth,sep=.3cm,center,rounded=true]{block body}
        \large \textbf{22/22}\\sorties SQL avec table exploitable
      \end{beamercolorbox}
    \column{.48\textwidth}
      \begin{beamercolorbox}[wd=\linewidth,sep=.3cm,center,rounded=true]{block body}
        \large \textbf{0}\\erreur SQL interne observée
      \end{beamercolorbox}
  \end{columns}
  \vfill
  {\scriptsize Validation technique du POC ; les scénarios ne constituent pas un échantillon d'usage en production.}
\end{frame}

% Include the report's original curve when the figure assets are available.
\IfFileExists{Figures/latency.png}{%
  \begin{frame}{Latence -- P50 et P99}
    \centering
    \reportfigure{Figures/latency.png}
      {Courbe P50/P99 de latence des appels LLM.}
      {Latence des appels LLM observée dans LangSmith -- figure du rapport.}
      {}
  \end{frame}
}{}

% 0:30
\begin{frame}{Résultat clé -- latence observée}
  \vspace{.75cm}
  \begin{columns}[T,onlytextwidth]
    \column{.32\textwidth}
      \kpi{3,01 s}{latence médiane}
    \column{.32\textwidth}
      \kpi{4,53 s}{latence moyenne}
    \column{.32\textwidth}
      \kpi{10,56 s}{latence p95}
  \end{columns}
  \vspace{.65cm}
  \begin{center}
    {\large Un maximum ponctuel de 71,52 s a été observé sur un scénario data complexe.}
  \end{center}
\end{frame}

% Include the report's original error chart when the figure assets are available.
\IfFileExists{Figures/error.png}{%
  \begin{frame}{Stabilité technique -- traces LangSmith}
    \centering
    \reportfigure{Figures/error.png}
      {Taux d'erreur observé sur les traces LangSmith.}
      {Taux d'erreur des traces pendant les évaluations du Supervisor -- figure du rapport.}
      {}
  \end{frame}
}{}

% 0:30
\begin{frame}{Discussion -- acquis et limites}
  \vspace{.35cm}
  \begin{columns}[T,onlytextwidth]
    \column{.47\textwidth}
      {\Large\bfseries Acquis}\\[.25cm]
      \begin{itemize}
        \item Socle data et backend mieux structurés
        \item Parcours métier traçables
        \item Architecture agentique modulaire
        \item Stabilité technique complète sur le jeu testé
      \end{itemize}
    \column{.47\textwidth}
      {\Large\bfseries Limites}\\[.25cm]
      \begin{itemize}
        \item Reprise de contexte multi-tours à renforcer
        \item Latence élevée sur certains cas ponctuels
        \item Évaluation et droits du RAG à approfondir
        \item Mesures du POC non représentatives de la production
      \end{itemize}
  \end{columns}
  \vfill
  \begin{center}
    \small Le gain manuel estimé dans le rapport est indicatif, non chronométré auprès d'utilisateurs.
  \end{center}
\end{frame}

% 0:30
\begin{frame}{Perspectives}
  \vspace{.45cm}
  \begin{itemize}
    \item \textbf{Fiabiliser} : compléter les tests multi-tours et persister les conversations
    \vspace{.38cm}
    \item \textbf{Améliorer le RAG} : recherche hybride, reranking, extraction PDF/OCR et évaluation des réponses
    \vspace{.38cm}
    \item \textbf{Renforcer la gouvernance} : droits d'accès, audit et sources documentaires autorisées
    \vspace{.38cm}
    \item \textbf{Industrialiser} : préparer Azure, une base PostgreSQL et la connexion aux sources réelles
  \end{itemize}
\end{frame}

% 0:25
\begin{frame}[plain]
  \vspace*{1.1cm}
  \begin{beamercolorbox}[wd=\paperwidth,sep=.5cm,center]{frametitle}
    {\color{white}\fontsize{28}{32}\selectfont\bfseries Conclusion}
  \end{beamercolorbox}
  \vspace{.9cm}
  \begin{center}
    {\LARGE\bfseries Un assistant IA en entreprise ne se résume pas à un prompt.}\\[.55cm]
    {\large Il repose sur des données structurées, des workflows contrôlés,
    des sorties exploitables et des garde-fous.}
  \end{center}
  \vfill
  \begin{center}
    {\color{StagePurple}\Huge\bfseries Merci pour votre attention}\\[.25cm]
    {\Large Questions}
  \end{center}
  \vfill
\end{frame}

\end{document}
